\section{\personFloer{}模空间的相容定向}
\label{sec:coherent-orientations-floer-moduli}

先从\personWitten{}与\personFloer{}的\personMorse{}同调观点，讨论经典\personMorse{}理论中的定向。

\subsection{\personMorse{}同调中的相容定向}
\label{subsec:coherent-orientations-morse-homology}

经典\personMorse{}理论中，负梯度流
$\dot x=-\operatorname{grad}f(x)$ 给出同调，而正梯度流给出上同调，这是熟知的事实
（见 (Mil65)）。考虑（向下）\personMorse{}复形 $CM(f)$ 与 $CM(-f)$，有自然同构
\[
 CM(-f)\cong\operatorname{Hom}(CM(f),\mathbb Q),
\]
它由一一对应
\begin{equation}
 \widetilde x\longmapsto x^*
 \label{eq:morse-dual-generator-correspondence}
\end{equation}
诱导，其中 $x^*\in\operatorname{Hom}(CM(f),\mathbb Q)$ 由
$x^*(x)=1$ 及对 $x\neq y\in\operatorname{Crit}(f)$ 有 $x^*(y)=0$ 定义。
这在\personMorse{}同调（更准确地说，\personMorse{}--\personSmale{}--\personWitten{}--\personFloer{}同调）的链层诱导
（强）非退化配对
\[
 CM_k(f)\times CM_{n-k}(-f)\longrightarrow\mathbb Q.
\]

按照赋予相容定向的标准方案(Mil65)，我们
\begin{enumerate}[label=(\arabic*)]
\item 先对每个 $x\in\operatorname{Crit}f$ 选取不稳定流形 $W_f^u(x)$ 的\emph{定向}；
\item 再选取稳定流形 $W_f^s(x)$ 的\emph{余定向}，使得作为有向向量空间有
\[
 N_xW_f^s(x)=T_xW_f^u(x).
\]
\end{enumerate}
\stfhIndex{I017}{00-ox}{$o(x)$}
把 $W_f^u(x)$ 的这个定向记为 $o(x)$。

\stfhIndex{I134}{Moersi-Simeierdui}{莫尔斯--斯梅尔对（Morse--Smale pair）}
对通有度量 $g$，稳定流形与不稳定流形横截相交，即 $(f,g)$ 成为\emph{\personMorse{}--\personSmale{}对}。
在该横截性假设下，要求
\[
 \operatorname{span}\{\nabla f(p)\}\oplus T_pS W_f^u(x)
 =T_pW_f^u(x)
\]
作为有向向量空间成立，由此赋予 $S W_f^u(x)$ 定向。由于
$W_f^s(y)\subset M$ 与 $S W_f^s(y)\subset f^{-1}(y)$ 在 $p$ 处的法丛典范同构，
把 $S W_f^s(y)$ 在 $f^{-1}(y)$ 中的余定向取为 $W_f^s(y)$ 在 $M$ 中的余定向。
\footnote{译者注：这里的 $y$ 同时表示临界点，原书却写层集 $f^{-1}(y)$；
紧接着又把该层写成 $f^{-1}(c)$。译文在本句照录原符号。}
因此，相交
\[
 S W_f^u(x)\cap S W_f^s(y)\subset f^{-1}(y)
\]
带有自然定向。并有自然微分同胚
\[
 S W_f^u(x)\cap S W_f^s(y)
 \cong\widetilde{\mathcal M}_g(f;x,y)/\mathbb R,
\]
其中 $\widetilde{\mathcal M}_g(f;x,y)$ 是\emph{负}梯度流线的模空间，即方程
\[
 \dot\chi+\operatorname{grad}f(\chi)=0,
 \qquad \chi(-\infty)=x,
 \quad \chi(\infty)=y
\]
\stfhIndex{I018}{00-oxy}{$o(x,y)$}
的解空间。把这个定向记为 $o(x,y)$。

\stfhSubIndex{I102}{Zhanhe}{粘合（gluing）}{Zhanheduanlieliuxian}{粘合断裂流线（gluing broken trajectories）}
这种赋予定向的方法在粘合构造下具有如下相容性。注意，
$\mathcal M_g(f;x,z)\times\mathcal M_g(f;z,y)$ 带有来自 $o(x,z)$ 与 $o(z,y)$
的自然乘积定向。考虑\termGluingMap{}
\[
 \operatorname{Glue}:\mathcal M_g(f;x,z)\times\mathcal M_g(f;z,y)
 \times[R,\infty)\longrightarrow\mathcal M_g(f;x,y),
\]
它光滑延拓为
\[
 \operatorname{Glue}:\mathcal M_g(f;x,z)\times\mathcal M_g(f;z,y)
 \times([R,\infty)\cup\{\infty\})\longrightarrow
 \overline{\mathcal M}_g(f;x,y).
\]
把乘积定向在 $\operatorname{Glue}$ 下推前到
$\mathcal M_g(f;x,z)\#\mathcal M_g(f;z,y)$ 所得的定向记为
\[
 o(x,z)\#o(z,y),
\]
\stfhIndex{I099}{Zhanhedingxiang}{粘合定向（glued orientation）}
称为\emph{粘合定向}。另一方面，
$\mathcal M_g(f;x,z)\#\mathcal M_g(f;z,y)$ 作为带边流形
\stfhIndex{I045}{Bianjiedingxiang}{边界定向（boundary orientation）}
$\overline{\mathcal M}_g(f;x,y)$ 的边界分支还带有另一定向，称为\emph{边界定向}，记作
$\partial o(x,y)$。

\begin{proposition}
\label{prop:morse-coherent-orientation-gluing}
以 $\overline{\mathcal M}_g(f;x,y)$ 表示
$\mathcal M_g(f;x,y)$ 的紧化，并假设在 $R=\infty$ 处限制\termGluingMap{}后，
\[
 \mathcal M_g(f;x,z)\times\mathcal M_g(f;z,y)
 \hookrightarrow\partial\overline{\mathcal M}_g(f;x,y)
\]
成为一个边界分支。则对所有满足
$\mathcal M(x,y;f)\neq\varnothing$ 的 $x,y\in\operatorname{Crit}f$，有
\begin{equation}
 o(x)=o(x,y)\#o(y).
 \label{eq:morse-orientation-factorization}
\end{equation}
特别地，
\begin{equation}
 o(x,z)\#o(z,y)=\partial o(x,y).
 \label{eq:morse-coherent-boundary-orientation}
\end{equation}
\end{proposition}

这里不证明该命题，只启发式地说明
\eqref{eq:morse-coherent-boundary-orientation}如何由
\eqref{eq:morse-orientation-factorization}推出。后者可形式地写为
$o(x,y)=o(x)-o(y)$；把 $o(x,z)$ 与 $o(z,y)$ 相加即得前者。

\stfhIndex{I057}{Xiangrongdingxiang}{相容定向（coherent orientation）}
一般地，把满足~\eqref{eq:morse-coherent-boundary-orientation}的定向系
$\{o(x,y)\}$ 称为所有 $x,y\in\operatorname{Crit}(f)$ 的\personMorse{}模空间
$\widetilde{\mathcal M}(x,y;f)$ 上的相容定向。上述构造给出一种几何地提供这种
相容系统的方法。把从定向\emph{不稳定}流形开始的程序所生成的相容定向系记作 $\sigma$。

当 $\mu_f(x)-\mu_f(y)=1$ 时，在 $f(y)<c<f(x)$ 的层级 $c$ 考虑相应稳定球与不稳定球
\[
 S W_f^u(x):=f^{-1}(c)\cap W_f^u(x),
 \qquad
 S W_f^s(y):=f^{-1}(c)\cap W_f^s(y).
\]
一般地，对与 $f^{-1}(c)$ 横截相交的子流形 $T\subset M$，相交
$T\cap f^{-1}(c)$ 的定向由关系
\begin{equation}
 T_pT=\{-\nabla f(p)\}\oplus T_p(T\cap f^{-1}(c))
 =\left\{\frac\partial{\partial\tau}\right\}
 \oplus T_p(T\cap f^{-1}(c))
 \label{eq:morse-level-intersection-orientation}
\end{equation}
给出。
\footnote{译者注：原书在本式左端写 $T_pK$，但本段定义的子流形为 $T$；
译文依本段对象补正。}
比较每个 $p\in S W_f^u(x)\cap S W_f^s(y)$ 处
$T_p(S W_f^u(x))$ 与 $N_p(S W_f^s(y))$ 上的两个定向，即可定义
$f^{-1}(c)$ 中的交数
$\#(S W_f^u(x)\cap S W_f^s(y))$。于是可以整数形式定义\personMorse{}同调边界映射
\[
 \partial_{(g,f)}:CM_*(f)\longrightarrow CM_*(f)
\]
的矩阵元
$n(f;x,y)=\#(\widetilde{\mathcal M}_g(f;x,y))$，并满足边界性质
$\partial_{(g,f)}^2=0$。相容系统 $\sigma$ 的相容性是建立该边界性质的关键。
定义关于定向相容系统 $\sigma$ 的 $(g,f)$ 的\personMorse{}同调为
\[
 HM_*^\sigma(g,f;M).
\]

显然，在上述构造中可用稳定流形取代不稳定流形，得到另一个相容定向系，记为
$\widetilde\sigma$。对该系统应注意：
\begin{enumerate}[label=(\arabic*)]
\item 先定向稳定流形；
\item 要求 $N^{\widetilde\sigma}W_f^u(x)=T^{\widetilde\sigma}W_f^s(x)$；
\item 对任意与 $M^c=f^{-1}(c)$ 横截相交的有向子流形 $S\subset M$，以关系
\[
 \left\{-\frac\partial{\partial\tau}\right\}
 \oplus T_p(S\cap M^c)=T_pS
\]
赋予 $S\cap M^c$ 定向。
\end{enumerate}
借助时间反演对合
\[
 \mathcal M(x,y;f)\longrightarrow\mathcal M(y,x;-f),
\]
可以把对 $f$ 构造 $\widetilde\sigma$ 视为对 $-f$ 构造 $\sigma$。一般而言
\[
 HM_k^{\widetilde\sigma}(g,f;M)\neq HM_k^\sigma(g,f;M),
\]
但稍后将在第 19.7.2 节证明，不论 $M$ 是否可定向，均有同构
\[
 HM_{n-k}^{\widetilde\sigma}(g,-f;M)\cong HM_k^\sigma(g,f;M).
\]
当 $M$ 可定向时，还可证明
\[
 HM_k^{\widetilde\sigma}(g,f;M)\cong HM_k^\sigma(g,f;M).
\]

\subsection{光滑\personFloer{}模空间上的定向数据}
\label{subsec:orientation-data-smooth-floer-moduli}

本节研究给定拉格朗日对 $(L_0,L_1)$ 的各个光滑\personFloer{}模空间的可定向性。
事实表明，关于\personMorse{}模空间定向的几何讨论，可以替换为对非线性映射
\begin{equation}
 \chi\longmapsto\dot\chi+\operatorname{grad}_gf(\chi)
 \label{eq:morse-nonlinear-gradient-map}
\end{equation}
之线性化\personFredholm{}算子的族指标所作的相应解析研究，即使在\personMorse{}同调语境中也是如此。
把~\eqref{eq:morse-nonlinear-gradient-map}换成 $\overline\partial$ 型非线性映射，即可在\personFloer{}理论语境中直接重复这一解析研究。

回忆，\personFloer{}同调理论常被视为环路空间（闭弦情形）或带拉格朗日边界条件的道路空间
（开弦情形）上的无穷维\personMorse{}理论。给\personFloer{}模空间赋予定向，可以视为给这一环路空间
或道路空间的一个子集赋予定向。众所周知，$M$ 的 spin 结构可以视为 $M$ 的环路空间的
一个定向。事实表明，要在相应的紧化模空间上提供相容定向，需要的是拉格朗日子流形
$L$ 的 spin 结构。
\footnote{译者注：原书此句写 ``spin structure of Lagrangian submanifold $M$''；
$M$ 在全章中是环境辛流形，按上下文这里应为拉格朗日子流形 $L$。译文据此补正。}
更一般地，需要如下定义。

\begin{definition}
\label{def:relatively-spin-lagrangian-chain}
若子流形 $L\subset M$ 可定向，并存在类
$\mathrm{st}\in H^2(M,\mathbb Z_2)$，使
\[
 \mathrm{st}|_L=w_2(TL),
\]
其中 $w_2(TL)$ 是 $TL$ 的\personStiefel{}--\personWhitney{}类，则称 $L$ 为\termRelativelySpin{}。

若存在类 $\mathrm{st}\in H^2(M,\mathbb Z_2)$，使对每个
$i=0,1,\ldots,k$ 均有
$\mathrm{st}|_{L_i}=w_2(TL_i)$，则称拉格朗日子流形链
$(L_0,L_1,\ldots,L_k)$ 或拉格朗日对 $(L_0,L_1)$ 为相对 spin。
\end{definition}

若 $N$ 是 spin 流形，则由嵌入 $\phi:N\to M$ 给出的任意拉格朗日子流形
$L=\phi(N)\subset(M,\omega)$ 都是相对 spin 的，因为可以取
$H^2(M,\mathbb Z_2)$ 中的 $\mathrm{st}=0$。

\stfhSubIndex{I058}{Xiangrongdingxiang}{相容定向（coherent orientation）}{ZaiFuluoertongdiaozhong}{在弗洛尔同调中（in Floer homology）}
\subsection{\personFloer{}同调中的相容定向}
\label{subsec:coherent-orientations-floer-homology}

下面说明一个相容定向系的构造，它对应于上一小节所述\personMorse{}同调的相容定向。
为简化叙述，只限于 spin 对 $(L_0,L_1)$。（关于\termRelativeSpinStructure{}的一般情形，
详见 (FOOO09) 第 8.1 节。）假设本节讨论中出现的所有拉格朗日子流形均配备定向与
spin 结构，并假设 $L_0$ 与 $L_1$ 横截相交。

本小节证明如下定理。

\begin{theorem}
\label{thm:spin-pair-orients-floer-moduli}
若 $(L_0,L_1)$ 是 spin 对，则 $\mathcal M(p,q;B)$ 可定向。当 $L_0,L_1$
为 spin 且选定 spin 结构后，对每个 $p,q\in L_0\cap L_1$ 及
$B\in\pi_2(p,q)$，它给出 $\mathcal M(p,q;B)$ 上的自然定向。
\end{theorem}

\subsubsection*{定向 $o_{[p,w]}$}
设 $[\ell_p,w]\in\Omega(L_0,L_1;\ell_{01})$，即 $p\in L_0\cap L_1$，且 $w$ 满足
\[
 w(0,t)=\ell_{01}(t),
 \quad w(1,t)\equiv p,
 \quad w(s,0)\in L_0,
 \quad w(s,1)\in L_1.
\]
\footnote{译者注：原书在此把 $[\widehat p,w]$ 与道路空间之间写成集合包含符号
$\subset$；译文按“一个元素”的文字说明改为 $\in$。}
以 $\operatorname{Map}(\ell_{01};p;L_0,L_1;\alpha)$ 表示此类映射
$w:[0,1]^2\to M$ 中满足
$[w]=\alpha\in\pi_2(\ell_{01};p)$ 者之集。取
$w\in\operatorname{Map}(\ell_{01};p;L_0,L_1;\alpha)$。同前，令
\[
 \Phi:w^*TM\longrightarrow[0,1]^2\times T_pM
\]
为一个（在同伦意义下唯一的）辛平凡化。平凡化 $\Phi$、边界条件
$w(0,t)=\ell_{01}(t)$ 以及沿 $\ell_{01}$ 的拉格朗日道路 $\lambda_{01}$，共同定义
一条拉格朗日道路
\[
 \lambda_\Phi=\lambda^\Phi_{([p,w];\lambda_{01})}:
 [0,1]\longrightarrow\Lambda^{\mathrm{ori}}(T_pM),
\]
满足 $\lambda_\Phi(0)=T_pL_0$、$\lambda_\Phi(1)=T_pL_1$。该道路的同伦类
不依赖平凡化 $\Phi$，只依赖 $[p,w]$ 及沿 $\ell_{01}$ 的 $\lambda_{01}$ 的同伦类。
与\personMorse{}同调不同，每个交点 $p\in L_0\cap L_1$ 都带有映射同伦类集
$\pi_2(\ell_{01};p)$。需要系统地同时选取各种
$\alpha\in\pi_2(\ell_{01};p)$ 所对应的~\eqref{eq:determinant-line-map-space}
之定向。

为此，把 $T_{\ell_{01}(0)}L_0$ 与 $T_{\ell_{01}(1)}L_1$ 上的给定定向延拓为
沿 $\ell_{01}$ 的拉格朗日道路 $\lambda_{01}$ 上的标架。考虑
$\mathbb R_{\geq0}\times[0,1]$ 上 $w^*TM$ 的 $W^{1,p}$ 级截面 $\xi$ 的边值问题
\begin{equation}
 \begin{cases}
 D_w\overline\partial(\xi)=0,\\
 \xi(0,t)\in\lambda_{01}(t),\\
 \xi(\tau,0)\in T_pL_0,
 \quad\xi(\tau,1)\in T_pL_1.
 \end{cases}
 \label{eq:half-strip-orientation-problem}
\end{equation}
这里 $D_w\overline\partial$ 是\personCauchy{}--\personRiemann{}方程的线性化算子。

记 $Z_-=\mathbb R_{\geq0}\times[0,1]$，并以
$W^{1,p}(Z_-;\lambda_{01})$ 表示 $Z_-$ 上满足上述边界条件的 $w^*TM$ 的
$W^{1,p}$ 截面集。于是得到\personFredholm{}算子
\[
 D_w\overline\partial:W^{1,p}(Z_-,T_pM;\lambda)
 \longrightarrow L^p(Z_-,T_pM\otimes\Lambda^{0,1}),
\]
\footnote{译者注：原书刚把定义域的边界条件记作 $\lambda_{01}$，此式却改用
未在本小节定义的 $\lambda$；下句算子记号又恢复为 $\lambda_{01}$。译文保留原式。}
记为 $\overline\partial_{([p,w];\lambda_{01})}$。其作为虚向量空间的指标为
\[
 \operatorname{Index}D_w\overline\partial
 =\ker D_w\overline\partial-\operatorname{coker}D_w\overline\partial,
\]
数值指标则由\personMaslov{}--\personMorse{}指标给出。把行列式线记作
$\det\overline\partial_{([p,w];\lambda_{01})}$。当 $w$ 在同伦类
$\alpha\in\pi_2(\ell_{01};p)=\pi_2(\ell_{01};p;L_0,L_1)$ 中变化时，
这些线定义线丛
\begin{equation}
 \det\overline\partial_{([p,w];\lambda_{01})}
 \longrightarrow\operatorname{Map}(\ell_{01};p;L_0,L_1;\alpha).
 \label{eq:determinant-line-map-space}
\end{equation}

下述命题在命题~\ref{prop:trivial-boundary-orients-disc-index} 的证明中已被隐式证明。

\begin{proposition}
\label{prop:determinant-line-trivial-spin-pair}
对一对有向 spin 拉格朗日子流形 $(L_0,L_1)$，
式~\eqref{eq:determinant-line-map-space}给出的线丛关于每个固定的
$\alpha$ 均平凡。
\end{proposition}

\begin{exercise}
\label{ex:extract-spin-triviality-proof}
从命题~\ref{prop:trivial-boundary-orients-disc-index}的证明中提取该命题的证明。
\end{exercise}

把线丛~\eqref{eq:determinant-line-map-space}的一个定向记为 $o_{[p,w]}$。
这个定向可看作\personMorse{}同调中临界点不稳定流形定向的对应物。如前所述，
不稳定流形有许多由 $\pi_2(p;\ell_{01})$ 参数化的连通分支。

令 $\Phi,\lambda_\Phi$ 如上。回忆，真正进入 $o_{[p,w]}$ 构造的数据是：
\begin{enumerate}[label=(\arabic*)]
\item 连接 $T_pL_0$ 与 $T_pL_1$ 的道路
$\lambda_\Phi:[0,1]\to\Lambda^{\mathrm{ori}}(T_pM)$；在定义 $o_{[p,w]}$ 时，该道路是在平凡化
$\Phi:w^*TM\to[0,1]^2\times T_pM$ 中的拼接
\[
 T_pL_0\longrightarrow T_{\ell_{01}(0)}L_0
 \#\lambda_{01}\#T_{\ell_{01}(1)}L_1\longrightarrow T_pL_1,
\]
连同 $w^*TM$ 在其上的相应拉格朗日子丛；
\footnote{译者注：原书此处把 $\lambda_\Phi$ 的目标写成 $T_pM$；但
$T_pL_0,T_pL_1$ 是拉格朗日子空间，且本页前文已把同一道路定义为取值于
$\Lambda^{\mathrm{ori}}(T_pM)$。译文据此前定义补正目标。}
\item 半带 $\mathbb R_{[0,\infty)}\times[0,1]$ 上、以该道路为边界条件的\personRiemann{}--\personHilbert{}问题~\eqref{eq:half-strip-orientation-problem}。
\end{enumerate}
因此需要用典范方式处理 $w$ 上平凡化 $\Phi$ 的选择。这是 $L_0,L_1$ 上 spin 结构
首次介入之处。回忆，$L_0,L_1$ 的定向分别决定
$T_pL_0=\lambda_{01}(0)$ 与 $T_pL_1=\lambda_{01}(1)$ 的定向。

\subsubsection*{定向 $o_p$}
考虑投影 $\operatorname{pr}:[0,1]\times M\to M$ 的拉回丛
$\operatorname{pr}^*TM\to[0,1]\times M$。对 $p\in L_0\cap L_1$，
取有向拉格朗日子空间道路
\[
 \lambda_p:[0,1]\longrightarrow\Lambda^{\mathrm{ori}}(T_pM),
 \qquad
 \lambda_p(0)=T_pL_0,
 \quad\lambda_p(1)=T_pL_1.
\]
滥用记号，也把连接 $T_pL_0,T_pL_1$ 的道路 $\lambda_p$ 视作丛
\[
 \lambda_p=\bigcup_{t\in[0,1]}\{t\}\times\lambda_p(t)\longrightarrow[0,1].
\]
该丛与 $i=0,1$ 时的 $[0,1]\times T_pL_i$ 一起，在
\[
 K(L_0,L_1)=(L_0\times\{0\})
 \cup((L_0\cap L_1)\times[0,1])
 \cup(L_1\times\{1\})\subset[0,1]\times M
\]
上定义 $\operatorname{pr}^*TM|_{K(L_0,L_1)}$ 的拉格朗日子丛
$\mathcal L\to K(L_0,L_1)$，它延拓子丛
\[
 \{0\}\times TL_0\cup\{1\}\times TL_1
 \subset\operatorname{pr}^*TM|_{\partial[0,1]}.
\]

以 $P_{\mathrm{spin}}(L_i)$ 表示 $TL_i$ 的主 spin 丛；它是有向标架丛
$P_{SO}(L_i)$ 的逐纤维双覆盖。可以假设 $L_0\cap L_1$ 包含在 $M$ 的 $3$-骨架中。
为在 $K(L_0,L_1)$ 的 $2$-骨架 $K(L_0,L_1)^{[2]}$ 上指定
$\mathcal L$ 的 spin 结构，作如下处理。

令 $\sigma$ 是 $P_{SO(n)}(\mathcal L)$ 的截面，即向量丛
$\mathcal L\to K(L_0,L_1)$ 的有向标架。考虑其在
$[0,1]\times\{p\}\subset K(L_0,L_1)$ 上的限制，记为
$\sigma_{[0,1]\times\{p\}}$。假设相伴定向对每个 $i=0,1$ 都限制为
$T_pL_i$ 的给定定向。回忆
$P_{Spin(n)}(\mathcal L)\to P_{SO(n)}(\mathcal L)$ 是双覆盖；对每个 $i=0,1$，令
\[
 \iota_{\sigma_i}(p)\in P_{\mathrm{spin}}(L_i)|_p
\]
是 $\sigma(i\times\{p\})$ 的一个提升。
\footnote{译者注：原书对两个 $i$ 都写成提升 $\sigma(1\times\{p\})$；
译文按两个端点及紧接着的条件~\eqref{eq:spin-lift-endpoint-compatibility}
补正为 $\sigma(i\times\{p\})$。}
考虑三元组
$(\sigma_{[0,1]\times\{p\}},\iota_{\sigma_0}(p),\iota_{\sigma_1}(p))$。
由于
\[
 P_{Spin(n)}(\mathcal L)|_{[0,1]\times\{p\}}
 \cong[0,1]\times Spin(n)
\]
是
$P_{SO(n)}(\mathcal L)|_{[0,1]\times\{p\}}
\cong[0,1]\times SO(n)$ 的平凡双覆盖，可以选取
$\sigma_{[0,1]\times\{p\}}$ 的提升
\[
 \iota_{[0,1]\times\{p\}}:[0,1]\times\{p\}
 \longrightarrow P_{Spin(n)}(\mathcal L)|_{[0,1]\times\{p\}},
\]
满足
\begin{equation}
 \iota_{[0,1]\times\{p\}}(i)=\iota_{\sigma_i}(p),
 \qquad i=0,1.
 \label{eq:spin-lift-endpoint-compatibility}
\end{equation}
这就在 $[0,1]\times\{p\}$ 上给出覆盖 $P_{SO(n)}(\mathcal L)$ 的截面
$\sigma$ 的 $\mathcal L$ 之 spin 结构。借助相容条件
\eqref{eq:spin-lift-endpoint-compatibility}，可以粘合 spin 丛
$P_{\mathrm{spin}}(\mathcal L\oplus V)|_{[0,1]\times\{p\}}$ 与
$P_{\mathrm{spin}}(\{i\}\times L_i)$（$i=0,1$），得到
$\mathcal L\to K(L_0,L_1)^{[2]}$ 的 spin 结构。
\footnote{译者注：原书在这里首次出现稳定化向量丛 $V$，但本段未定义 $V$。
译文照录该符号。}

随后，对这条典范道路 $\alpha_p$，利用
$\operatorname{pr}^*TM\to K(L_0,L_1)$ 的拉格朗日子丛
$\mathcal L\to K(L_0,L_1)$ 上上述固定的 spin 结构，重复 $o_{[p,w]}$ 的构造，
得到\termDeterminantLineBundle{}
\[
 \det\overline\partial_{([p,w];\alpha_p)}
\]
的定向。对每个 $p\in L_0\cap L_1$，把这一方向记为 $o_p$。

\subsubsection*{$\mathcal M(p,q;B)$ 上的相容定向}
回忆，$\overline{\mathcal M}(p,r;B)$ 的边界
$\partial\overline{\mathcal M}(p,r;B)$ 是下列各连通分支之并：
\[
 \mathcal M(p,q;B_1)\#\mathcal M(q,r;B_2),
 \qquad B=B_1\#B_2.
\]
\footnote{译者注：源书印 $\partial\mathcal M(p,r;B)$（无杠），译文按断裂模空间边界补杠，
记作 $\partial\overline{\mathcal M}(p,r;B)$。}
以 $\partial o(p,r;B)$ 表示边界
$\partial\overline{\mathcal M}(p,r;B)$ 的诱导边界定向。另一方面，以
\[
 o(p,q;B_1)\#o(q,r;B_2)
\]
表示由粘合\termVirtualTangentSpace{}自然得到的乘积
$\mathcal M(p,q;B_1)\#\mathcal M(q,r;B_2)$ 上的自然定向。

\begin{definition}[相容定向系]
\label{def:coherent-system-floer-orientations}
若对 $p,q\in L_0\cap L_1$ 及 $B\in\pi_2(p,q)$ 的定向族
$\{o(p,q;B)\}$，在
$\mathcal M(p,q;B_1)\#\mathcal M(q,r;B_2)
\subset\partial\overline{\mathcal M}(p,r;B)$ 且
$\mathcal M(p,r;B)$ 的虚维数为 $1$ 时，满足粘合规则
\begin{equation}
 \partial o(p,r;B)=o(p,q;B_1)\#o(q,r;B_2),
 \label{eq:coherent-floer-orientation-gluing-rule}
\end{equation}
则称它为相容定向系。
\end{definition}

下面是\personFloer{}同调中命题~\ref{prop:morse-coherent-orientation-gluing}的对应物。

\begin{theorem}
\label{thm:construction-coherent-floer-orientations}
设 $(L_0,L_1)$ 是一对横截相交的有向 spin 拉格朗日子流形。对每个
$p\in L_0\cap L_1$，令 $\alpha_p$ 为
式~\eqref{eq:reference-lagrangian-path-alpha-p}给出的 $T_pM$ 中的拉格朗日道路。
对所有 $p$ 固定 $\operatorname{Index}\overline\partial_{\alpha_p}$ 上定向 $o_p$ 的选择。
\begin{enumerate}[label=(\arabic*)]
\item 它决定~\eqref{eq:determinant-line-map-space}上的定向，记为 $o_{[p,w]}$。
\item 此外，$o_p,o_{[p,w]}$ 通过粘合规则
\begin{equation}
 o_{[q,w\#B]}=o_{[p,w]}\#o(p,q;B)
 \label{eq:orientation-capping-gluing-rule}
\end{equation}
决定所有 $p,q\in L_0\cap L_1$、$B\in\pi_2(p,q)$ 的
$\mathcal M(p,q;B)$ 上定向 $o(p,q;B)$。
\item 如此定义的系统 $\{o(p,q;B)\}$ 是相容的。
\end{enumerate}
\end{theorem}

\begin{proof}
由命题~\ref{prop:determinant-line-trivial-spin-pair}，线丛
\eqref{eq:determinant-line-map-space}平凡。对每个 $p$，给定椭圆算子
\[
 \overline\partial_{\alpha_p,Z_+}:W_{\alpha_p}^{1,p}(Z_+;T_pM)
 \longrightarrow L^p(Z_+;T_pM\otimes\Lambda^{0,1}(Z_+)).
\]
对 $R>0$，如下定义空间 $(\mathbb R\times[0,1])\#_RZ_+$。考虑
\[
 Z_{+,R}=\{z\in Z_-\mid\operatorname{Re}z\geq-R\}.
\]
\footnote{译者注：本章前面把 $Z_-$ 定义为 $\operatorname{Re}z\geq0$、
$Z_+$ 定义为 $\operatorname{Re}z\leq0$；因此这里的
$Z_{+,R}\subset Z_-$ 既与下标不一致，条件 $\operatorname{Re}z\geq-R$
在 $Z_-$ 上也不产生所需截断。译文照录原式。}
粘合 $[0,R]\times[0,1]$ 与 $Z_{+,R}$：分别把
$(R,t)\in Z_{-,R}$ 与 $(-R,t)\in[-R,R]\times[0,1]$ 等同，并把
$(R,t)\in[-R,R]\times[0,1]$ 与 $(-R,t)\in Z_{+,R}$ 等同。
\footnote{译者注：本句宣称只粘合两个空间，却使用三个记号
$Z_{-,R}$、$[-R,R]\times[0,1]$、$Z_{+,R}$，且前一个未定义；译文照录其
等同关系，不替原书重构定义域。}

为证 (1)，以显然方式“粘合”算子
$D_w\overline\partial_{(\mathbb R_{\geq0}\times[0,1])}$ 与
问题~\eqref{eq:half-strip-orientation-problem}，得到
$(\mathbb R_{\geq0}\times[0,1])\#Z_+$ 上的算子
$D_w\overline\partial_{(\mathbb R_{\geq0}\times[0,1])\#Z_+}$。
有（算子族的）虚向量空间同构
\begin{equation}
 \operatorname{Index}
 \bigl(D_w\overline\partial_{(\mathbb R_{\geq0}\times[0,1])\#Z_+}\bigr)
 \cong\operatorname{Index}\overline\partial_{([p,w];\lambda_{01})}
 \oplus\operatorname{Index}\overline\partial_{\alpha_p}.
 \label{eq:orientation-index-gluing-isomorphism}
\end{equation}
需要指出，空间
$\mathbb R_{\geq0}\times[0,1]\#_RZ_+$ 共形于一个二维圆盘；其上的拉回切丛
在边界上带有具有 spin 结构的拉格朗日子丛。

我们固定了 $\alpha_p$ 与 $\lambda_{01}$ 的正向标架，它们延拓由 $L_0,L_1$
上的 spin 结构在其二骨架上提供的 $TL_0,TL_1$ 标架（或平凡化）。由
命题~\ref{prop:determinant-line-trivial-spin-pair}，该平凡化诱导指标丛
$\operatorname{Index}
(D_w\overline\partial_{(\mathbb R_{\geq0}\times[0,1])\#Z_+})$
的典范定向。因此，$\operatorname{Index}\overline\partial_{\alpha_p}$ 的定向
（连同该典范定向）典范地诱导
$\operatorname{Index}\overline\partial_{([p,w];\lambda_{01})}$ 的定向。这证明 (1)。

为证 (2)，与上面类似地把 $Z_-$ 与 $\mathbb R\times[0,1]$ 粘合，即得结论，
细节从略。最后，(3) 是~\eqref{eq:orientation-capping-gluing-rule}的推论。
证明完成。
\end{proof}

顺便指出，定理~\ref{thm:construction-coherent-floer-orientations}
也蕴含定理~\ref{thm:spin-pair-orients-floer-moduli}。
